\chapter{Dynamic Programming: DP on Subsequences (Knapsack Family)}
\chapterepigraph{Almost every problem in this chapter is the same
question asked three different ways: does a subset with property $X$
exist, how many subsets have property $X$, or what is the best value
achievable subject to property $X$. Settle which of the three you are
answering before writing a single line of recurrence.}

\section{The Knapsack Family, and Why It All Starts with Subset Sum}
\label{sec:dp-subseq-intro}

The Recursion chapter's Power Set section's pick/not-pick backtracking tree
generates every subsequence of an array in $O(2^n)$. This chapter asks
DP's defining question about that same tree: \emph{do many different
pick/not-pick paths reach the same $(\textit{index}, \textit{target})$
state}, and if so, can we cache by that state to collapse the
exponential tree into a polynomial DP table? For every problem in this
chapter, the answer is yes, and the state is almost always
$(\textit{index}, \textit{remainingTarget})$ -- the same two parameters
the Recursion chapter's Combination Sum section already recursed on, now with a
memo table layered underneath.

\begin{patternbox}[The Three Questions, and the 0/1 vs.\ Unbounded Distinction]
\begin{itemize}
  \item \textbf{``Does a subset summing to $K$ exist?''} -- a
        \textbf{boolean} DP table, $\textit{dp}[\textit{index}][\textit{target}]$,
        combining the pick and not-pick branches with a logical OR.
  \item \textbf{``How many subsets sum to $K$?''} -- the same table,
        but combining pick and not-pick with \textbf{addition} (their
        counts are disjoint and exhaustive, exactly as in
        the Recursion chapter's Count Subsequences with Sum K section).
  \item \textbf{``What is the maximum value achievable, subject to a
        capacity constraint?''} -- the same table shape, combining pick
        and not-pick with \textbf{maximum} instead of OR or addition --
        this is the 0/1 Knapsack problem, the family's namesake.
\end{itemize}
A second, independent axis determines whether the \textbf{pick} branch
advances to $\textit{index}+1$ (each item usable \textbf{once} -- ``0/1''
knapsack, as in the Recursion chapter's Subset Sum II section) or stays at
$\textit{index}$ (each item usable \textbf{unlimited} times --
``unbounded'' knapsack, as in the Recursion chapter's Combination Sum section).
This chapter covers the 0/1 family first (Sections~\ref{sec:dp-subset-sum-target}--\ref{sec:dp-01-knapsack})
before the unbounded family
(Sections~\ref{sec:dp-coin-change}--\ref{sec:dp-rod-cutting}).
\end{patternbox}

Every problem here is carried through the same four stages as every
earlier DP chapter -- recursion, memoization, tabulation, space
optimization -- with one new wrinkle in the tabulation stage: because the
state has \emph{two} dimensions (index and target/capacity), the table
is genuinely 2D, and space optimization (when the recurrence only reads
the immediately previous index's row) reduces it to a single rolling 1D
array, exactly as grid DP's rolling-row trick did in the previous
chapter -- just with ``target/capacity'' playing the role ``column''
played there.

\newpage

\section{Subset Sum Equal to Target}
\label{sec:dp-subset-sum-target}

\problemstatement{Given an array $\textit{arr}$ of $n$ non-negative
integers and a target $k$, determine whether \textbf{some} subset of
$\textit{arr}$ sums to exactly $k$.}

\begin{patternbox}
This is the Recursion chapter's ``does a
subsequence with sum $K$ exist'' recursion, revisited now that we know
to look for overlapping subproblems: the state
$(\textit{index}, \textit{remaining})$ is reached along many different
pick/not-pick paths, so caching by that pair of parameters converts the
$O(2^n)$ recursion into an $O(n \cdot k)$ DP.
\end{patternbox}

\subsection*{Deriving the recurrence}

Let $f(\textit{index}, \textit{target})$ be \texttt{true} if some subset
of $\textit{arr}[0..\textit{index}]$ sums to exactly $\textit{target}$.
At index $\textit{index}$, either this element is excluded
($f(\textit{index}-1, \textit{target})$ must hold), or it is included
(requiring $\textit{arr}[\textit{index}] \le \textit{target}$, and
$f(\textit{index}-1, \textit{target}-\textit{arr}[\textit{index}])$ must
hold):
\[
  f(\textit{index}, \textit{target}) = f(\textit{index}{-}1, \textit{target})
  \ \lor\ \big(\textit{arr}[\textit{index}] \le \textit{target} \ \land\
  f(\textit{index}{-}1, \textit{target}-\textit{arr}[\textit{index}])\big).
\]
Base case: $f(0, \textit{target}) = (\textit{target}=0) \lor
(\textit{target} = \textit{arr}[0])$.

\subsection*{Approach 1: Recursion}

\begin{lstlisting}
#include <bits/stdc++.h>
using namespace std;

bool subsetSumRecursive(int index, int target, vector<int>& arr) {
    if (target == 0) return true;
    if (index == 0) return arr[0] == target;

    bool notPick = subsetSumRecursive(index - 1, target, arr);
    bool pick = false;
    if (arr[index] <= target) {
        pick = subsetSumRecursive(index - 1, target - arr[index], arr);
    }
    return notPick || pick;
}

int main() {
    vector<int> arr = {1, 2, 3, 4};
    int n = arr.size();
    cout << (subsetSumRecursive(n - 1, 6, arr) ? "true" : "false") << endl;
    return 0;
}
\end{lstlisting}

\begin{complexitybox}[Approach 1 Complexity]
\textbf{Time.} $O(2^n)$.
\textbf{Space.} $O(n)$ recursion stack.
\end{complexitybox}

\subsection*{Approach 2: Memoization}

\begin{lstlisting}
#include <bits/stdc++.h>
using namespace std;

bool subsetSumMemo(int index, int target, vector<int>& arr,
                    vector<vector<int>>& memo) {
    if (target == 0) return true;
    if (index == 0) return arr[0] == target;
    if (memo[index][target] != -1) return memo[index][target];

    bool notPick = subsetSumMemo(index - 1, target, arr, memo);
    bool pick = false;
    if (arr[index] <= target) {
        pick = subsetSumMemo(index - 1, target - arr[index], arr, memo);
    }
    return memo[index][target] = (notPick || pick);
}

int main() {
    vector<int> arr = {1, 2, 3, 4};
    int n = arr.size(), k = 6;
    vector<vector<int>> memo(n, vector<int>(k + 1, -1));
    cout << (subsetSumMemo(n - 1, k, arr, memo) ? "true" : "false") << endl;
    return 0;
}
\end{lstlisting}

\begin{complexitybox}[Approach 2 Complexity]
\textbf{Time.} $O(n \cdot k)$.
\textbf{Space.} $O(n \cdot k)$ memo $+$ $O(n)$ recursion stack.
\end{complexitybox}

\subsection*{Approach 3: Tabulation}

\begin{lstlisting}
#include <bits/stdc++.h>
using namespace std;

bool subsetSumTabulation(vector<int>& arr, int k) {
    int n = arr.size();
    vector<vector<bool>> dp(n, vector<bool>(k + 1, false));

    for (int index = 0; index < n; index++) dp[index][0] = true;
    if (arr[0] <= k) dp[0][arr[0]] = true;

    for (int index = 1; index < n; index++) {
        for (int target = 1; target <= k; target++) {
            bool notPick = dp[index - 1][target];
            bool pick = false;
            if (arr[index] <= target) {
                pick = dp[index - 1][target - arr[index]];
            }
            dp[index][target] = notPick || pick;
        }
    }
    return dp[n - 1][k];
}

int main() {
    vector<int> arr = {1, 2, 3, 4};
    cout << (subsetSumTabulation(arr, 6) ? "true" : "false") << endl;
    return 0;
}
\end{lstlisting}

\subsection*{Dry run of the tabulation table}

\dryrun{Take $\textit{arr} = [1,2,3,4]$, $k=6$.}

\begin{center}
\begin{tabular}{c|ccccccc}
$\textit{index}\backslash\textit{target}$ & 0 & 1 & 2 & 3 & 4 & 5 & 6 \\ \hline
$0$ (val $1$) & T & T & F & F & F & F & F \\
$1$ (val $2$) & T & T & T & T & F & F & F \\
$2$ (val $3$) & T & T & T & T & T & T & T \\
$3$ (val $4$) & T & T & T & T & T & T & T
\end{tabular}
\end{center}

At $\textit{index}=2$ ($\textit{arr}[2]=3$), $\textit{target}=6$:
not-pick gives $\textit{dp}[1][6]=\text{F}$; pick gives
$\textit{dp}[1][3]=\text{T}$ (since $3\le6$). So
$\textit{dp}[2][6]=\text{T}$, already confirming a subset exists
($\{2,4\}$ sums to $6$, found even before considering index $3$). The
final answer, $\textit{dp}[3][6]=\text{T}$, confirms it.

\begin{complexitybox}[Approach 3 Complexity]
\textbf{Time.} $O(n \cdot k)$.
\textbf{Space.} $O(n \cdot k)$ table, no recursion stack.
\end{complexitybox}

\subsection*{Approach 4: Space Optimization}

\begin{lstlisting}
#include <bits/stdc++.h>
using namespace std;

bool subsetSumSpaceOptimized(vector<int>& arr, int k) {
    int n = arr.size();
    vector<bool> prev(k + 1, false);

    prev[0] = true;
    if (arr[0] <= k) prev[arr[0]] = true;

    for (int index = 1; index < n; index++) {
        vector<bool> current(k + 1, false);
        current[0] = true;
        for (int target = 1; target <= k; target++) {
            bool notPick = prev[target];
            bool pick = (arr[index] <= target) ? prev[target - arr[index]] : false;
            current[target] = notPick || pick;
        }
        prev = current;
    }
    return prev[k];
}

int main() {
    vector<int> arr = {1, 2, 3, 4};
    cout << (subsetSumSpaceOptimized(arr, 6) ? "true" : "false") << endl;
    return 0;
}
\end{lstlisting}

\begin{complexitybox}[Approach 4 Complexity]
\textbf{Time.} $O(n \cdot k)$.
\textbf{Space.} $O(k)$ for the rolling row, instead of $O(n \cdot k)$.
\end{complexitybox}

\begin{takeawaybox}
This boolean-OR recurrence, indexed by $(\textit{index},
\textit{target})$, is the foundation of every 0/1-knapsack-family
problem in this chapter -- Partition Equal Subset Sum, Count Subsets
with Sum K, and Count Partitions with a Given Difference are all this
\emph{exact} table, either queried at a derived target or with the OR
replaced by addition for counting. Internalize this one table fully
before moving on, since nothing new is derived from here through
Section~\ref{sec:dp-count-partitions-diff}.
\end{takeawaybox}

\section{Partition Equal Subset Sum}
\label{sec:dp-partition-equal-subset}

\problemstatement{Given an array $\textit{arr}$, determine whether it can
be partitioned into \textbf{two} subsets with \textbf{equal} sum.}

\begin{patternbox}
\emph{``Partition into two equal-sum subsets''} reduces directly to
Section~\ref{sec:dp-subset-sum-target}: if the whole array can be split
into two subsets of equal sum $S$, then the total sum must be $2S$ (so
an \textbf{odd} total sum makes this immediately impossible), and finding
one subset summing to exactly $S = \textit{totalSum}/2$ automatically
gives the other subset too (everything not in the first subset).
\end{patternbox}

\subsection*{Understanding the reduction}

No new recurrence needs deriving: compute $\textit{totalSum}$; if it is
odd, return \texttt{false} immediately. Otherwise, ask exactly
Section~\ref{sec:dp-subset-sum-target}'s question with target
$\textit{totalSum}/2$. If a subset summing to $\textit{totalSum}/2$
exists, the remaining elements necessarily sum to the same value
(since the two parts must add up to $\textit{totalSum} =
2 \times \textit{totalSum}/2$), giving a valid equal partition.

\subsection*{All Four Approaches}

Every one of the four stages is Section~\ref{sec:dp-subset-sum-target}'s
implementation, called with $k = \textit{totalSum}/2$. Only the
space-optimized version is shown here, to avoid repeating the identical
code four times.

\begin{lstlisting}
#include <bits/stdc++.h>
using namespace std;

bool subsetSumSpaceOptimized(vector<int>& arr, int k) {
    int n = arr.size();
    vector<bool> prev(k + 1, false);

    prev[0] = true;
    if (arr[0] <= k) prev[arr[0]] = true;

    for (int index = 1; index < n; index++) {
        vector<bool> current(k + 1, false);
        current[0] = true;
        for (int target = 1; target <= k; target++) {
            bool notPick = prev[target];
            bool pick = (arr[index] <= target) ? prev[target - arr[index]] : false;
            current[target] = notPick || pick;
        }
        prev = current;
    }
    return prev[k];
}

bool canPartition(vector<int>& arr) {
    int totalSum = accumulate(arr.begin(), arr.end(), 0);
    if (totalSum % 2 != 0) return false;

    return subsetSumSpaceOptimized(arr, totalSum / 2);
}

int main() {
    vector<int> arr = {1, 5, 11, 5};
    cout << (canPartition(arr) ? "true" : "false") << endl; // true
    return 0;
}
\end{lstlisting}

\subsection*{Dry run}

\dryrun{Take $\textit{arr} = [1,5,11,5]$.}

$\textit{totalSum} = 22$, even, so target $= 11$. Running
Section~\ref{sec:dp-subset-sum-target}'s subset-sum check for target
$11$: the subset $\{5,5,1\}$ sums to $11$ -- found, so the answer is
\texttt{true}. The complementary subset $\{11\}$ also sums to $11$,
confirming the equal partition $\{1,5,5\}$ and $\{11\}$.

\begin{complexitybox}
\textbf{Time Complexity.} $O(n \cdot \textit{totalSum})$, inherited
directly from Section~\ref{sec:dp-subset-sum-target}'s subset-sum DP
with $k = \textit{totalSum}/2$.
\textbf{Space Complexity.} $O(\textit{totalSum})$ for the rolling row
(space-optimized version).
\end{complexitybox}

\begin{takeawaybox}
Before deriving a new DP recurrence, always check whether the problem is
a thin wrapper around one already solved -- here, an odd-sum early exit
plus a single call to subset-sum with a derived target. This is the same
``reduce to an already-solved sub-problem'' instinct already exercised
for House Robber II's circular-array reduction back in
the DP 1D chapter.
\end{takeawaybox}

\section{Count Subsets with Sum K}
\label{sec:dp-count-subsets-sum-k}

\problemstatement{Given an array $\textit{arr}$ of non-negative integers
and a target $k$, count \textbf{how many} subsets sum to exactly $k$.}

\begin{patternbox}
This is Section~\ref{sec:dp-subset-sum-target}'s exact recurrence with
the logical OR replaced by \textbf{addition} -- the same
exists-to-count transformation as the Recursion chapter's
``Check if There Exists a Subsequence with Sum K'' to ``Count
Subsequences with Sum K'' transition,
now applied on top of the memoized/tabulated table instead of plain
recursion.
\end{patternbox}

\subsection*{Deriving the recurrence}

\[
  f(\textit{index}, \textit{target}) = f(\textit{index}{-}1, \textit{target})
  \ +\ \big(\textit{arr}[\textit{index}] \le \textit{target}\ ?\
  f(\textit{index}{-}1, \textit{target}-\textit{arr}[\textit{index}]) : 0\big).
\]
Base case (handling a possible zero-valued element carefully, since a
$0$ can be included or excluded without changing the sum, doubling the
count): if $\textit{index}=0$, return $2$ if $\textit{target}=0$ and
$\textit{arr}[0]=0$; return $1$ if $\textit{target}=0$ or
$\textit{target}=\textit{arr}[0]$ (but not both conditions above); else
return $0$.

\subsection*{Approach 1: Recursion}

\begin{lstlisting}
#include <bits/stdc++.h>
using namespace std;

int countSubsetsRecursive(int index, int target, vector<int>& arr) {
    if (index == 0) {
        if (target == 0 && arr[0] == 0) return 2;
        if (target == 0 || target == arr[0]) return 1;
        return 0;
    }

    int notPick = countSubsetsRecursive(index - 1, target, arr);
    int pick = 0;
    if (arr[index] <= target) {
        pick = countSubsetsRecursive(index - 1, target - arr[index], arr);
    }
    return notPick + pick;
}

int main() {
    vector<int> arr = {1, 2, 2, 3};
    int n = arr.size();
    cout << countSubsetsRecursive(n - 1, 3, arr) << endl; // 3
    return 0;
}
\end{lstlisting}

\begin{complexitybox}[Approach 1 Complexity]
\textbf{Time.} $O(2^n)$.
\textbf{Space.} $O(n)$ recursion stack.
\end{complexitybox}

\subsection*{Approach 2: Memoization}

\begin{lstlisting}
#include <bits/stdc++.h>
using namespace std;

int countSubsetsMemo(int index, int target, vector<int>& arr,
                      vector<vector<int>>& memo) {
    if (index == 0) {
        if (target == 0 && arr[0] == 0) return 2;
        if (target == 0 || target == arr[0]) return 1;
        return 0;
    }
    if (memo[index][target] != -1) return memo[index][target];

    int notPick = countSubsetsMemo(index - 1, target, arr, memo);
    int pick = 0;
    if (arr[index] <= target) {
        pick = countSubsetsMemo(index - 1, target - arr[index], arr, memo);
    }
    return memo[index][target] = notPick + pick;
}

int main() {
    vector<int> arr = {1, 2, 2, 3};
    int n = arr.size(), k = 3;
    vector<vector<int>> memo(n, vector<int>(k + 1, -1));
    cout << countSubsetsMemo(n - 1, k, arr, memo) << endl; // 3
    return 0;
}
\end{lstlisting}

\begin{complexitybox}[Approach 2 Complexity]
\textbf{Time.} $O(n \cdot k)$.
\textbf{Space.} $O(n \cdot k)$ memo $+$ $O(n)$ recursion stack.
\end{complexitybox}

\subsection*{Approach 3: Tabulation}

\begin{lstlisting}
#include <bits/stdc++.h>
using namespace std;

int countSubsetsTabulation(vector<int>& arr, int k) {
    int n = arr.size();
    vector<vector<int>> dp(n, vector<int>(k + 1, 0));

    dp[0][0] = (arr[0] == 0) ? 2 : 1;
    if (arr[0] != 0 && arr[0] <= k) dp[0][arr[0]] = 1;

    for (int index = 1; index < n; index++) {
        for (int target = 0; target <= k; target++) {
            int notPick = dp[index - 1][target];
            int pick = (arr[index] <= target) ? dp[index - 1][target - arr[index]] : 0;
            dp[index][target] = notPick + pick;
        }
    }
    return dp[n - 1][k];
}

int main() {
    vector<int> arr = {1, 2, 2, 3};
    cout << countSubsetsTabulation(arr, 3) << endl; // 3
    return 0;
}
\end{lstlisting}

\subsection*{Dry run of the tabulation table}

\dryrun{Take $\textit{arr} = [1,2,2,3]$, $k=3$.}

\begin{center}
\begin{tabular}{c|c|c|c|c}
$\textit{index}\backslash\textit{target}$ & 0 & 1 & 2 & 3 \\ \hline
$0$ (val $1$) & 1 & 1 & 0 & 0 \\
$1$ (val $2$) & 1 & 1 & 1 & 1 \\
$2$ (val $2$) & 1 & 1 & 2 & 2 \\
$3$ (val $3$) & 1 & 1 & 2 & 3
\end{tabular}
\end{center}

The three qualifying subsets summing to $3$ are $\{1,2\}$ (using the
first $2$), $\{1,2\}$ (using the second $2$), and $\{3\}$ alone --
matching $\textit{dp}[3][3]=3$.

\begin{complexitybox}[Approach 3 Complexity]
\textbf{Time.} $O(n \cdot k)$.
\textbf{Space.} $O(n \cdot k)$ table, no recursion stack.
\end{complexitybox}

\subsection*{Approach 4: Space Optimization}

\begin{lstlisting}
#include <bits/stdc++.h>
using namespace std;

int countSubsetsSpaceOptimized(vector<int>& arr, int k) {
    int n = arr.size();
    vector<int> prev(k + 1, 0);

    prev[0] = (arr[0] == 0) ? 2 : 1;
    if (arr[0] != 0 && arr[0] <= k) prev[arr[0]] = 1;

    for (int index = 1; index < n; index++) {
        vector<int> current(k + 1, 0);
        for (int target = 0; target <= k; target++) {
            int notPick = prev[target];
            int pick = (arr[index] <= target) ? prev[target - arr[index]] : 0;
            current[target] = notPick + pick;
        }
        prev = current;
    }
    return prev[k];
}

int main() {
    vector<int> arr = {1, 2, 2, 3};
    cout << countSubsetsSpaceOptimized(arr, 3) << endl; // 3
    return 0;
}
\end{lstlisting}

\begin{complexitybox}[Approach 4 Complexity]
\textbf{Time.} $O(n \cdot k)$.
\textbf{Space.} $O(k)$ for the rolling row.
\end{complexitybox}

\begin{takeawaybox}
Counting subsets and checking subset existence differ by exactly one
operator (addition instead of logical OR) at every step of an otherwise
identical recurrence -- the same lesson already drawn from the
Recursion chapter's exists-to-count transition.
Watch for the zero-value edge case specifically whenever a problem allows
zero-valued array elements, since a $0$ can always be freely included or
excluded without affecting the sum, silently doubling some counts if not
handled in the base case.
\end{takeawaybox}

\section{Count Partitions with a Given Difference}
\label{sec:dp-count-partitions-diff}

\problemstatement{Given an array $\textit{arr}$ and an integer $d$,
count the number of ways to partition $\textit{arr}$ into two subsets
$S_1$ and $S_2$ such that $S_1 - S_2 = d$ (where $S_1, S_2$ denote the
sums of the two subsets, and every element belongs to exactly one of
them).}

\begin{patternbox}
Another reduction to an already-solved problem: given
$S_1 - S_2 = d$ and $S_1 + S_2 = \textit{totalSum}$ (every element is in
exactly one subset), solving this small system of two equations gives
$S_1 = (\textit{totalSum}+d)/2$ -- turning the partition-counting
question into exactly Section~\ref{sec:dp-count-subsets-sum-k}'s
``count subsets summing to $K$,'' with $K = (\textit{totalSum}+d)/2$.
\end{patternbox}

\subsection*{Understanding the reduction}

Adding the two equations $S_1-S_2=d$ and $S_1+S_2=\textit{totalSum}$
eliminates $S_2$: $2S_1 = \textit{totalSum}+d$, so $S_1 =
(\textit{totalSum}+d)/2$. If $\textit{totalSum}+d$ is odd, no integer
$S_1$ satisfies this, so the answer is immediately $0$. Otherwise, every
way of choosing a subset summing to exactly this $S_1$ corresponds to
exactly one valid partition (the chosen subset is $S_1$, everything else
is $S_2$, and $S_2 = \textit{totalSum}-S_1$ automatically satisfies
$S_1-S_2=d$) -- so the answer is exactly
Section~\ref{sec:dp-count-subsets-sum-k}'s subset-count, evaluated at
this derived target.

\subsection*{All Four Approaches}

Every stage reuses Section~\ref{sec:dp-count-subsets-sum-k}'s
implementation unchanged, called with $k = (\textit{totalSum}+d)/2$.
Only the space-optimized version, plus the reduction wrapper, is shown.

\begin{lstlisting}
#include <bits/stdc++.h>
using namespace std;

int countSubsetsSpaceOptimized(vector<int>& arr, int k) {
    int n = arr.size();
    vector<int> prev(k + 1, 0);

    prev[0] = (arr[0] == 0) ? 2 : 1;
    if (arr[0] != 0 && arr[0] <= k) prev[arr[0]] = 1;

    for (int index = 1; index < n; index++) {
        vector<int> current(k + 1, 0);
        for (int target = 0; target <= k; target++) {
            int notPick = prev[target];
            int pick = (arr[index] <= target) ? prev[target - arr[index]] : 0;
            current[target] = notPick + pick;
        }
        prev = current;
    }
    return prev[k];
}

int countPartitionsWithDiff(vector<int>& arr, int d) {
    int totalSum = accumulate(arr.begin(), arr.end(), 0);
    if ((totalSum + d) % 2 != 0 || totalSum + d < 0) return 0;

    int s1 = (totalSum + d) / 2;
    return countSubsetsSpaceOptimized(arr, s1);
}

int main() {
    vector<int> arr = {1, 1, 2, 3};
    cout << countPartitionsWithDiff(arr, 1) << endl; // 3
    return 0;
}
\end{lstlisting}

\subsection*{Dry run}

\dryrun{Take $\textit{arr} = [1,1,2,3]$, $d=1$.}

$\textit{totalSum}=7$. $S_1 = (7+1)/2 = 4$. Counting subsets summing to
$4$: $\{1,3\}$ (using the first $1$), $\{1,3\}$ (using the second $1$),
and $\{1,1,2\}$ -- three subsets, giving $\textit{answer}=3$. Checking
one of them: subset $\{1,3\}$ has sum $4$; the complementary subset
$\{1,2\}$ has sum $3$; difference $4-3=1=d$, confirmed.

\begin{complexitybox}
\textbf{Time Complexity.} $O(n \cdot \textit{totalSum})$, inherited
from Section~\ref{sec:dp-count-subsets-sum-k}.
\textbf{Space Complexity.} $O(\textit{totalSum})$ for the rolling row.
\end{complexitybox}

\begin{takeawaybox}
Whenever a problem describes a partition in terms of a \textbf{sum and a
difference} of the two parts, solving the resulting two-equation system
for one part's target sum is almost always the fastest path to reusing
an already-solved subset-counting DP, rather than inventing a new
two-dimensional recurrence that tracks both parts directly.
\end{takeawaybox}

\section{A Note on Assign Cookies: When DP Is the Wrong Tool}
\label{sec:dp-assign-cookies-note}

\noindent Striver's sheet lists Assign Cookies again here, inside the DP
on Subsequences lecture -- and it is worth pausing on \emph{why}, since
the honest answer is a valuable lesson in its own right: \textbf{this
problem does not need dynamic programming at all}, and forcing a DP
solution onto it would be strictly worse than the greedy solution
already derived in the Greedy chapter.

\begin{patternbox}[Recognising When a Subset/Assignment Problem Is Actually Greedy]
Assign Cookies asks: given greed factors $g$ and cookie sizes $s$,
maximize the number of children satisfiable, where child $i$ is
satisfied by cookie $j$ if $s[j] \ge g[i]$. This \emph{looks} like an
assignment problem that backtracking or DP might solve via some
subset-matching recurrence -- but the Greedy chapter's Assign Cookies section
already proved, via a direct exchange argument, that sorting both arrays
and greedily matching the smallest sufficient cookie to the smallest
unsatisfied child is \textbf{provably optimal}. No DP table is needed
because there is no genuine trade-off between overlapping sub-problems
to resolve: the greedy choice at each step is never wrong, so there is
nothing for memoization to usefully cache.
\end{patternbox}

\subsection*{Why a DP/backtracking solution would be strictly worse here}

A backtracking solution exploring every possible assignment of cookies to
children would cost at least $O(2^{n})$ in the worst case (each child
either gets a cookie or doesn't, combinatorially), or, with a more
clever subset-based DP state, perhaps $O(n \cdot m)$ or worse depending
on formulation -- but the Greedy chapter's greedy algorithm
already solves the problem in $O(n \log n + m \log m)$ (dominated by
sorting) with $O(1)$ extra space. Since the greedy algorithm is already
proven optimal, any DP formulation can only ever \emph{match} this
bound at best, typically with strictly worse constants and memory
usage -- there is no scenario in which reaching for DP here is the
better engineering decision.

\begin{ideabox}[The Real Lesson]
Before reaching for memoization or tabulation on a subset/assignment-flavored
problem, always ask first: \emph{can a sorted greedy pass already solve
this optimally?} The keyword-recognition habits from the Greedy chapter's
Assign Cookies section, in particular, should be
checked \emph{before} defaulting to a DP recurrence, not after. DP is
the right tool precisely when greedy's exchange argument \emph{fails} --
when the locally best choice can depend on choices made elsewhere, which
is exactly what does \emph{not} happen in Assign Cookies.
\end{ideabox}

\begin{takeawaybox}
Not every problem that involves choosing a subset or making an
assignment needs dynamic programming. Recognizing when a simpler,
already-optimal greedy algorithm applies -- and confidently \emph{not}
reaching for a DP table in that case -- is as important a skill as
knowing how to build the DP table when one is genuinely needed.
\end{takeawaybox}

\section{Target Sum}
\label{sec:dp-target-sum}

\problemstatement{Given an array $\textit{nums}$ and an integer
$\textit{target}$, assign either a $+$ or $-$ sign to each element, and
count the number of distinct sign assignments that make the total
expression evaluate to exactly $\textit{target}$.}

\begin{patternbox}
\emph{``Assign $+$ or $-$ to each element to reach a target sum''} is,
word for word, Section~\ref{sec:dp-count-partitions-diff}'s partition
problem in disguise: every element assigned $+$ belongs to one ``subset''
$P$ (positive sum), every element assigned $-$ belongs to the other
subset $N$ (negative sum), and the expression's value is exactly
$P - N$. Counting sign assignments reaching $\textit{target}$ is
therefore exactly counting partitions with difference
$d = \textit{target}$ -- \textbf{no new derivation is needed at all.}
\end{patternbox}

\subsection*{Understanding the reduction}

Every element of $\textit{nums}$ goes into exactly one of two groups:
those assigned $+$ (summing to $P$) and those assigned $-$ (summing to
$N$, as a positive magnitude). The expression's total value is $P-N$,
and $P+N$ is always the array's total sum (every element contributes its
magnitude to exactly one of the two groups). This is precisely
Section~\ref{sec:dp-count-partitions-diff}'s setup, with $d =
\textit{target}$: count subsets summing to $(\textit{totalSum} +
\textit{target})/2$.

\subsection*{All Four Approaches}

Every stage is Section~\ref{sec:dp-count-partitions-diff}'s
implementation, called unchanged with $d = \textit{target}$.

\begin{lstlisting}
#include <bits/stdc++.h>
using namespace std;

int countSubsetsSpaceOptimized(vector<int>& arr, int k) {
    int n = arr.size();
    vector<int> prev(k + 1, 0);

    prev[0] = (arr[0] == 0) ? 2 : 1;
    if (arr[0] != 0 && arr[0] <= k) prev[arr[0]] = 1;

    for (int index = 1; index < n; index++) {
        vector<int> current(k + 1, 0);
        for (int target = 0; target <= k; target++) {
            int notPick = prev[target];
            int pick = (arr[index] <= target) ? prev[target - arr[index]] : 0;
            current[target] = notPick + pick;
        }
        prev = current;
    }
    return prev[k];
}

int findTargetSumWays(vector<int>& nums, int target) {
    int totalSum = accumulate(nums.begin(), nums.end(), 0);
    if ((totalSum + target) % 2 != 0 || totalSum + target < 0) return 0;

    int p = (totalSum + target) / 2;
    return countSubsetsSpaceOptimized(nums, p);
}

int main() {
    vector<int> nums = {1, 1, 1, 1, 1};
    cout << findTargetSumWays(nums, 3) << endl; // 5
    return 0;
}
\end{lstlisting}

\subsection*{Dry run}

\dryrun{Take $\textit{nums} = [1,1,1,1,1]$, $\textit{target}=3$.}

$\textit{totalSum}=5$. $P = (5+3)/2 = 4$. Counting subsets of five
$1$s summing to $4$: there are $\binom{5}{4}=5$ ways to choose which
four of the five $1$s are positive (the remaining one is negative),
giving $5$ valid sign assignments -- each evaluates to
$1+1+1+1-1=3=\textit{target}$, confirming the answer.

\begin{complexitybox}
\textbf{Time Complexity.} $O(n \cdot \textit{totalSum})$, inherited
from Section~\ref{sec:dp-count-subsets-sum-k}.
\textbf{Space Complexity.} $O(\textit{totalSum})$ for the rolling row.
\end{complexitybox}

\begin{takeawaybox}
This section is intentionally almost content-free beyond the reduction
itself: once Count Partitions with a Given Difference is solved, Target
Sum requires recognizing the disguise and nothing more. A surprising
number of seemingly distinct DP problems on interview sheets are exactly
this -- the same two or three core recurrences, wearing different cover
stories.
\end{takeawaybox}

\section{0/1 Knapsack}
\label{sec:dp-01-knapsack}

\problemstatement{Given $n$ items, where item $i$ has weight
$\textit{wt}[i]$ and value $\textit{val}[i]$, and a knapsack of capacity
$W$, choose a subset of items (each usable \textbf{at most once}) whose
total weight does not exceed $W$, maximizing total value.}

\begin{patternbox}
This is the chapter's namesake, and it generalizes
Section~\ref{sec:dp-subset-sum-target}'s boolean reachability table in
the most direct way possible: instead of asking ``\emph{can} we reach
exactly capacity $\textit{target}$,'' we ask ``what is the
\emph{maximum value} achievable using \emph{at most} capacity $W$'' --
replacing the logical OR combination with a \textbf{maximum}, and adding
each picked item's \emph{value} (not just consuming its weight) to the
running total. Contrast this directly with
the Greedy chapter's Fractional Knapsack problem: there,
items could be split, which let a single greedy ratio-sort solve it
optimally in $O(n \log n)$; here, items are all-or-nothing, which
destroys that greedy argument and is exactly why this problem needs a
genuine DP table.
\end{patternbox}

\subsection*{Deriving the recurrence}

Let $f(\textit{index}, \textit{capacity})$ be the maximum value
achievable using items $0..\textit{index}$, subject to a weight budget
of $\textit{capacity}$. At item $\textit{index}$: either exclude it
($f(\textit{index}-1, \textit{capacity})$), or include it (only if
$\textit{wt}[\textit{index}] \le \textit{capacity}$, gaining
$\textit{val}[\textit{index}]$ plus
$f(\textit{index}-1, \textit{capacity}-\textit{wt}[\textit{index}])$) --
take whichever is larger:
\begin{gather*}
  f(\textit{index}, \textit{capacity}) = \max\Big(
  f(\textit{index}{-}1, \textit{capacity}),\ \textit{take} \Big), \\
  \textit{take} = \textit{wt}[\textit{index}] \le \textit{capacity}\ ?\
  \textit{val}[\textit{index}] + f(\textit{index}{-}1,
  \textit{capacity}-\textit{wt}[\textit{index}]) : -\infty.
\end{gather*}
Base case: $f(0, \textit{capacity}) = \textit{val}[0]$ if
$\textit{wt}[0] \le \textit{capacity}$, else $0$.

\subsection*{Approach 1: Recursion}

\begin{lstlisting}
#include <bits/stdc++.h>
using namespace std;

int knapsackRecursive(int index, int capacity, vector<int>& wt, vector<int>& val) {
    if (index == 0) {
        return (wt[0] <= capacity) ? val[0] : 0;
    }

    int notPick = knapsackRecursive(index - 1, capacity, wt, val);
    int pick = INT_MIN;
    if (wt[index] <= capacity) {
        pick = val[index] + knapsackRecursive(index - 1, capacity - wt[index], wt, val);
    }
    return max(notPick, pick);
}

int main() {
    vector<int> wt = {1, 3, 4, 5};
    vector<int> val = {1, 4, 5, 7};
    int n = wt.size(), W = 7;
    cout << knapsackRecursive(n - 1, W, wt, val) << endl; // 9
    return 0;
}
\end{lstlisting}

\begin{complexitybox}[Approach 1 Complexity]
\textbf{Time.} $O(2^n)$.
\textbf{Space.} $O(n)$ recursion stack.
\end{complexitybox}

\subsection*{Approach 2: Memoization}

\begin{lstlisting}
#include <bits/stdc++.h>
using namespace std;

int knapsackMemo(int index, int capacity, vector<int>& wt, vector<int>& val,
                  vector<vector<int>>& memo) {
    if (index == 0) {
        return (wt[0] <= capacity) ? val[0] : 0;
    }
    if (memo[index][capacity] != -1) return memo[index][capacity];

    int notPick = knapsackMemo(index - 1, capacity, wt, val, memo);
    int pick = INT_MIN;
    if (wt[index] <= capacity) {
        pick = val[index] + knapsackMemo(index - 1, capacity - wt[index], wt, val, memo);
    }
    return memo[index][capacity] = max(notPick, pick);
}

int main() {
    vector<int> wt = {1, 3, 4, 5};
    vector<int> val = {1, 4, 5, 7};
    int n = wt.size(), W = 7;
    vector<vector<int>> memo(n, vector<int>(W + 1, -1));
    cout << knapsackMemo(n - 1, W, wt, val, memo) << endl; // 9
    return 0;
}
\end{lstlisting}

\begin{complexitybox}[Approach 2 Complexity]
\textbf{Time.} $O(n \cdot W)$.
\textbf{Space.} $O(n \cdot W)$ memo $+$ $O(n)$ recursion stack.
\end{complexitybox}

\subsection*{Approach 3: Tabulation}

\begin{lstlisting}
#include <bits/stdc++.h>
using namespace std;

int knapsackTabulation(vector<int>& wt, vector<int>& val, int W) {
    int n = wt.size();
    vector<vector<int>> dp(n, vector<int>(W + 1, 0));

    for (int capacity = wt[0]; capacity <= W; capacity++) dp[0][capacity] = val[0];

    for (int index = 1; index < n; index++) {
        for (int capacity = 0; capacity <= W; capacity++) {
            int notPick = dp[index - 1][capacity];
            int pick = INT_MIN;
            if (wt[index] <= capacity) {
                pick = val[index] + dp[index - 1][capacity - wt[index]];
            }
            dp[index][capacity] = max(notPick, pick);
        }
    }
    return dp[n - 1][W];
}

int main() {
    vector<int> wt = {1, 3, 4, 5};
    vector<int> val = {1, 4, 5, 7};
    cout << knapsackTabulation(wt, val, 7) << endl; // 9
    return 0;
}
\end{lstlisting}

\subsection*{Dry run of the tabulation table}

\dryrun{Take $\textit{wt}=[1,3,4,5]$, $\textit{val}=[1,4,5,7]$, $W=7$.}

\begin{center}
\begin{tabular}{c|cccccccc}
$\textit{index}\backslash W$ & 0 & 1 & 2 & 3 & 4 & 5 & 6 & 7 \\ \hline
$0$ (wt$1$,val$1$) & 0 & 1 & 1 & 1 & 1 & 1 & 1 & 1 \\
$1$ (wt$3$,val$4$) & 0 & 1 & 1 & 4 & 5 & 5 & 5 & 5 \\
$2$ (wt$4$,val$5$) & 0 & 1 & 1 & 4 & 5 & 6 & 6 & 9 \\
$3$ (wt$5$,val$7$) & 0 & 1 & 1 & 4 & 5 & 7 & 8 & 9
\end{tabular}
\end{center}

At $\textit{index}=2$, $W=7$: not-pick gives $\textit{dp}[1][7]=5$; pick
gives $\textit{val}[2]+\textit{dp}[1][3] = 5+4=9$. Max is $9$. The final
answer is $\textit{dp}[3][7]=9$, achieved by items $1$ (wt$3$, val$4$)
and $2$ (wt$4$, val$5$) together: total weight $3+4=7\le7$, total value
$4+5=9$, matching the table.

\begin{complexitybox}[Approach 3 Complexity]
\textbf{Time.} $O(n \cdot W)$.
\textbf{Space.} $O(n \cdot W)$ table, no recursion stack.
\end{complexitybox}

\subsection*{Approach 4: Space Optimization}

\begin{lstlisting}
#include <bits/stdc++.h>
using namespace std;

int knapsackSpaceOptimized(vector<int>& wt, vector<int>& val, int W) {
    int n = wt.size();
    vector<int> prev(W + 1, 0);

    for (int capacity = wt[0]; capacity <= W; capacity++) prev[capacity] = val[0];

    for (int index = 1; index < n; index++) {
        vector<int> current(W + 1, 0);
        for (int capacity = 0; capacity <= W; capacity++) {
            int notPick = prev[capacity];
            int pick = INT_MIN;
            if (wt[index] <= capacity) {
                pick = val[index] + prev[capacity - wt[index]];
            }
            current[capacity] = max(notPick, pick);
        }
        prev = current;
    }
    return prev[W];
}

int main() {
    vector<int> wt = {1, 3, 4, 5};
    vector<int> val = {1, 4, 5, 7};
    cout << knapsackSpaceOptimized(wt, val, 7) << endl; // 9
    return 0;
}
\end{lstlisting}

\begin{complexitybox}[Approach 4 Complexity]
\textbf{Time.} $O(n \cdot W)$.
\textbf{Space.} $O(W)$ for the rolling row.
\end{complexitybox}

\begin{takeawaybox}
0/1 Knapsack is the most general member of this chapter's 0/1 family:
Subset Sum is the special case where every item's weight equals its
value (so reachability and value-maximization coincide), and Count
Subsets is the special case that counts qualifying selections instead of
optimizing a value. Recognizing 0/1 Knapsack as the ``parent'' recurrence
of everything from Section~\ref{sec:dp-subset-sum-target} onward
consolidates this entire half of the chapter into one shape.
\end{takeawaybox}

\section{Coin Change (Minimum Coins)}
\label{sec:dp-coin-change}

\noindent We now begin the chapter's \textbf{unbounded knapsack} family:
every problem from here on lets each item (coin denomination, rod
length) be used an \textbf{unlimited} number of times, which means the
pick branch stays at the same index instead of advancing -- exactly the
recursion shape the Recursion chapter's Combination Sum section
already derived.

\problemstatement{Given an array $\textit{coins}$ of denominations
(each available in \textbf{unlimited} supply) and a target
$\textit{amount}$, find the \textbf{minimum} number of coins needed to
make exactly $\textit{amount}$, or $-1$ if it is impossible.}

\begin{patternbox}
\emph{``Unlimited supply of each denomination, minimize the count
used''} is the unbounded-knapsack analogue of
Section~\ref{sec:dp-subset-sum-target}'s boolean reachability table: the
pick branch stays at the same index (since the same coin may be reused),
and the combination operator is a \textbf{minimum} of counts instead of a
logical OR of reachability.
\end{patternbox}

\subsection*{Deriving the recurrence}

Let $f(\textit{index}, \textit{amount})$ be the minimum number of coins
(using denominations $0..\textit{index}$) needed to make
$\textit{amount}$. Either we exclude denomination $\textit{index}$
entirely ($f(\textit{index}-1, \textit{amount})$), or we use it at least
once (paying $1$ coin, then recursing on the \emph{same} index since it
remains available, with the reduced amount):
\begin{gather*}
  f(\textit{index}, \textit{amount}) = \min\Big(
  f(\textit{index}{-}1, \textit{amount}),\ \textit{take} \Big), \\
  \textit{take} = \textit{coins}[\textit{index}] \le \textit{amount}\ ?\
  1 + f(\textit{index}, \textit{amount}-\textit{coins}[\textit{index}]) :
  \infty.
\end{gather*}
Base case: $f(0, \textit{amount}) = \textit{amount} /
\textit{coins}[0]$ if $\textit{amount}$ is exactly divisible by
$\textit{coins}[0]$, else $\infty$ (using only the first denomination,
either it divides evenly or it is impossible).

\subsection*{Approach 1: Recursion}

\begin{lstlisting}
#include <bits/stdc++.h>
using namespace std;
const int INF = 1e9;

int coinChangeRecursive(int index, int amount, vector<int>& coins) {
    if (index == 0) {
        return (amount % coins[0] == 0) ? amount / coins[0] : INF;
    }

    int notTake = coinChangeRecursive(index - 1, amount, coins);
    int take = INF;
    if (coins[index] <= amount) {
        take = 1 + coinChangeRecursive(index, amount - coins[index], coins);
    }
    return min(notTake, take);
}

int main() {
    vector<int> coins = {1, 2, 5};
    int n = coins.size(), amount = 11;
    int result = coinChangeRecursive(n - 1, amount, coins);
    cout << (result >= INF ? -1 : result) << endl; // 3
    return 0;
}
\end{lstlisting}

\begin{complexitybox}[Approach 1 Complexity]
\textbf{Time.} Exponential -- each call can recurse on the same index
many times (reducing the amount), giving a recursion tree that grows
combinatorially with $\textit{amount}$.
\textbf{Space.} $O(\textit{amount})$ recursion stack in the worst case
(repeatedly taking the smallest coin).
\end{complexitybox}

\subsection*{Approach 2: Memoization}

\begin{lstlisting}
#include <bits/stdc++.h>
using namespace std;
const int INF = 1e9;

int coinChangeMemo(int index, int amount, vector<int>& coins,
                    vector<vector<int>>& memo) {
    if (index == 0) {
        return (amount % coins[0] == 0) ? amount / coins[0] : INF;
    }
    if (memo[index][amount] != -1) return memo[index][amount];

    int notTake = coinChangeMemo(index - 1, amount, coins, memo);
    int take = INF;
    if (coins[index] <= amount) {
        take = 1 + coinChangeMemo(index, amount - coins[index], coins, memo);
    }
    return memo[index][amount] = min(notTake, take);
}

int main() {
    vector<int> coins = {1, 2, 5};
    int n = coins.size(), amount = 11;
    vector<vector<int>> memo(n, vector<int>(amount + 1, -1));
    int result = coinChangeMemo(n - 1, amount, coins, memo);
    cout << (result >= INF ? -1 : result) << endl; // 3
    return 0;
}
\end{lstlisting}

\begin{complexitybox}[Approach 2 Complexity]
\textbf{Time.} $O(n \cdot \textit{amount})$.
\textbf{Space.} $O(n \cdot \textit{amount})$ memo $+$
$O(\textit{amount})$ recursion stack.
\end{complexitybox}

\subsection*{Approach 3: Tabulation}

\begin{lstlisting}
#include <bits/stdc++.h>
using namespace std;
const int INF = 1e9;

int coinChangeTabulation(vector<int>& coins, int amount) {
    int n = coins.size();
    vector<vector<int>> dp(n, vector<int>(amount + 1, INF));

    for (int a = 0; a <= amount; a++) {
        if (a % coins[0] == 0) dp[0][a] = a / coins[0];
    }

    for (int index = 1; index < n; index++) {
        for (int a = 0; a <= amount; a++) {
            int notTake = dp[index - 1][a];
            int take = INF;
            if (coins[index] <= a) {
                take = 1 + dp[index][a - coins[index]];
            }
            dp[index][a] = min(notTake, take);
        }
    }
    int result = dp[n - 1][amount];
    return result >= INF ? -1 : result;
}

int main() {
    vector<int> coins = {1, 2, 5};
    cout << coinChangeTabulation(coins, 11) << endl; // 3
    return 0;
}
\end{lstlisting}

\subsection*{Dry run of the tabulation table}

\dryrun{Take $\textit{coins} = [1,2,5]$, $\textit{amount}=11$.}

Row $0$ (only coin $1$): $\textit{dp}[0][a] = a$ for all $a$ (every
amount needs exactly $a$ coins of value $1$). Row $1$ (coins $1,2$):
e.g.\ $\textit{dp}[1][11] = \min(\textit{dp}[0][11],\ 1+\textit{dp}[1][9])$;
unrolling, using coin $2$ as many times as possible gives
$\textit{dp}[1][11]=6$ ($5$ twos and one... actually $11=2\times5+1$,
so $5$ coins of $2$ plus $1$ coin of $1$ $=6$ coins). Row $2$ (coins
$1,2,5$): $\textit{dp}[2][11] = \min(\textit{dp}[1][11],\
1+\textit{dp}[2][6])$; following this through, the optimal is
$5+5+1=11$ using $3$ coins, giving $\textit{dp}[2][11]=3$.

\begin{complexitybox}[Approach 3 Complexity]
\textbf{Time.} $O(n \cdot \textit{amount})$.
\textbf{Space.} $O(n \cdot \textit{amount})$ table, no recursion stack.
\end{complexitybox}

\subsection*{Approach 4: Space Optimization}

\begin{lstlisting}
#include <bits/stdc++.h>
using namespace std;
const int INF = 1e9;

int coinChangeSpaceOptimized(vector<int>& coins, int amount) {
    int n = coins.size();
    vector<int> prev(amount + 1, INF);

    for (int a = 0; a <= amount; a++) {
        if (a % coins[0] == 0) prev[a] = a / coins[0];
    }

    for (int index = 1; index < n; index++) {
        vector<int> current(amount + 1, INF);
        for (int a = 0; a <= amount; a++) {
            int notTake = prev[a];
            int take = INF;
            if (coins[index] <= a) {
                take = 1 + current[a - coins[index]]; // same row! (unbounded)
            }
            current[a] = min(notTake, take);
        }
        prev = current;
    }
    int result = prev[amount];
    return result >= INF ? -1 : result;
}

int main() {
    vector<int> coins = {1, 2, 5};
    cout << coinChangeSpaceOptimized(coins, 11) << endl; // 3
    return 0;
}
\end{lstlisting}

\begin{ideabox}[Why ``Take'' Reads from \textit{current}, Not \textit{prev}]
Because the pick branch in an \textbf{unbounded} knapsack recurrence
stays at the \emph{same} index, the space-optimized version must read
the take-branch value from the row currently being built
($\textit{current}$), not from the previous row ($\textit{prev}$) --
the opposite of every 0/1 knapsack-family problem earlier in this
chapter, where both branches read only from $\textit{prev}$. This
single detail is the only change space optimization needs when moving
from the 0/1 family to the unbounded family.
\end{ideabox}

\begin{complexitybox}[Approach 4 Complexity]
\textbf{Time.} $O(n \cdot \textit{amount})$.
\textbf{Space.} $O(\textit{amount})$ for the rolling row.
\end{complexitybox}

\begin{takeawaybox}
Coin Change is 0/1 Knapsack's unbounded twin: same boolean-to-minimum
combination idea, but with the pick branch staying at the same index --
exactly the same 0/1-versus-unbounded distinction the Recursion
chapter drew between Subset Sum II (no reuse) and Combination Sum
(unlimited reuse). The single space-optimization wrinkle above (reading
``take'' from the current row) is the one new implementation detail this
unbounded family introduces.
\end{takeawaybox}

\section{Coin Change II (Count the Ways)}
\label{sec:dp-coin-change-2}

\problemstatement{Given an array $\textit{coins}$ of denominations (each
in unlimited supply) and a target $\textit{amount}$, count the
\textbf{number of distinct combinations} of coins that sum to exactly
$\textit{amount}$.}

\begin{patternbox}
This is Section~\ref{sec:dp-coin-change}'s recurrence with minimum
replaced by \textbf{addition} -- the same exists-vs-count distinction
drawn throughout this chapter and the Recursion chapter before it, now
applied to the unbounded-reuse setting.
\end{patternbox}

\subsection*{Deriving the recurrence}

\begin{gather*}
  f(\textit{index}, \textit{amount}) = f(\textit{index}{-}1, \textit{amount})
  \ +\ \textit{take}, \\
  \textit{take} = \textit{coins}[\textit{index}] \le \textit{amount}\ ?\
  f(\textit{index}, \textit{amount}-\textit{coins}[\textit{index}]) : 0.
\end{gather*}
Base case: $f(0, \textit{amount}) = 1$ if $\textit{amount}$ is exactly
divisible by $\textit{coins}[0]$ (exactly one way: use that many copies
of the first coin), else $0$.

\subsection*{Approach 1: Recursion}

\begin{lstlisting}
#include <bits/stdc++.h>
using namespace std;

int changeRecursive(int index, int amount, vector<int>& coins) {
    if (index == 0) {
        return (amount % coins[0] == 0) ? 1 : 0;
    }

    int notTake = changeRecursive(index - 1, amount, coins);
    int take = 0;
    if (coins[index] <= amount) {
        take = changeRecursive(index, amount - coins[index], coins);
    }
    return notTake + take;
}

int main() {
    vector<int> coins = {1, 2, 5};
    int n = coins.size();
    cout << changeRecursive(n - 1, 5, coins) << endl; // 4
    return 0;
}
\end{lstlisting}

\begin{complexitybox}[Approach 1 Complexity]
\textbf{Time.} Exponential, same overlapping-subproblem blowup as
Section~\ref{sec:dp-coin-change}.
\textbf{Space.} $O(\textit{amount})$ recursion stack.
\end{complexitybox}

\subsection*{Approach 2: Memoization}

\begin{lstlisting}
#include <bits/stdc++.h>
using namespace std;

int changeMemo(int index, int amount, vector<int>& coins,
                vector<vector<int>>& memo) {
    if (index == 0) {
        return (amount % coins[0] == 0) ? 1 : 0;
    }
    if (memo[index][amount] != -1) return memo[index][amount];

    int notTake = changeMemo(index - 1, amount, coins, memo);
    int take = 0;
    if (coins[index] <= amount) {
        take = changeMemo(index, amount - coins[index], coins, memo);
    }
    return memo[index][amount] = notTake + take;
}

int main() {
    vector<int> coins = {1, 2, 5};
    int n = coins.size(), amount = 5;
    vector<vector<int>> memo(n, vector<int>(amount + 1, -1));
    cout << changeMemo(n - 1, amount, coins, memo) << endl; // 4
    return 0;
}
\end{lstlisting}

\begin{complexitybox}[Approach 2 Complexity]
\textbf{Time.} $O(n \cdot \textit{amount})$.
\textbf{Space.} $O(n \cdot \textit{amount})$ memo $+$
$O(\textit{amount})$ recursion stack.
\end{complexitybox}

\subsection*{Approach 3: Tabulation}

\begin{lstlisting}
#include <bits/stdc++.h>
using namespace std;

int changeTabulation(vector<int>& coins, int amount) {
    int n = coins.size();
    vector<vector<long long>> dp(n, vector<long long>(amount + 1, 0));

    for (int a = 0; a <= amount; a++) {
        dp[0][a] = (a % coins[0] == 0) ? 1 : 0;
    }

    for (int index = 1; index < n; index++) {
        for (int a = 0; a <= amount; a++) {
            long long notTake = dp[index - 1][a];
            long long take = (coins[index] <= a) ? dp[index][a - coins[index]] : 0;
            dp[index][a] = notTake + take;
        }
    }
    return dp[n - 1][amount];
}

int main() {
    vector<int> coins = {1, 2, 5};
    cout << changeTabulation(coins, 5) << endl; // 4
    return 0;
}
\end{lstlisting}

\subsection*{Dry run of the tabulation table}

\dryrun{Take $\textit{coins}=[1,2,5]$, $\textit{amount}=5$.}

The four ways to make $5$: $\{1,1,1,1,1\}$, $\{1,1,1,2\}$,
$\{1,2,2\}$, $\{5\}$ -- matching $\textit{dp}[2][5]=4$. Row $0$ (coin
$1$ only): $\textit{dp}[0][a]=1$ for every $a$ (always exactly one way
using only $1$s). Row $1$ (coins $1,2$): combinations using some
$2$s and the rest $1$s -- $\textit{dp}[1][5]=3$ (zero, one, or two
$2$s, with the remainder in $1$s). Row $2$ adds the option of using
coin $5$ once, contributing the fourth way.

\begin{complexitybox}[Approach 3 Complexity]
\textbf{Time.} $O(n \cdot \textit{amount})$.
\textbf{Space.} $O(n \cdot \textit{amount})$ table, no recursion stack.
\end{complexitybox}

\subsection*{Approach 4: Space Optimization}

\begin{lstlisting}
#include <bits/stdc++.h>
using namespace std;

int changeSpaceOptimized(vector<int>& coins, int amount) {
    int n = coins.size();
    vector<long long> prev(amount + 1, 0);

    for (int a = 0; a <= amount; a++) {
        prev[a] = (a % coins[0] == 0) ? 1 : 0;
    }

    for (int index = 1; index < n; index++) {
        vector<long long> current(amount + 1, 0);
        for (int a = 0; a <= amount; a++) {
            long long notTake = prev[a];
            long long take = (coins[index] <= a) ? current[a - coins[index]] : 0;
            current[a] = notTake + take;
        }
        prev = current;
    }
    return prev[amount];
}

int main() {
    vector<int> coins = {1, 2, 5};
    cout << changeSpaceOptimized(coins, 5) << endl; // 4
    return 0;
}
\end{lstlisting}

\begin{complexitybox}[Approach 4 Complexity]
\textbf{Time.} $O(n \cdot \textit{amount})$.
\textbf{Space.} $O(\textit{amount})$ for the rolling row, reading
``take'' from \textit{current} exactly as in
Section~\ref{sec:dp-coin-change}.
\end{complexitybox}

\begin{takeawaybox}
Coin Change and Coin Change II are, structurally, the exact same
recurrence with minimum versus addition as the only difference -- the
unbounded-knapsack equivalent of Subset Sum versus Count Subsets earlier
in this chapter. Once you have one of a min/count/exists triplet for a
given recurrence shape, the other two cost almost nothing extra to
derive.
\end{takeawaybox}

\section{Rod Cutting}
\label{sec:dp-rod-cutting}

\problemstatement{Given a rod of length $n$ and an array $\textit{price}$
where $\textit{price}[i]$ is the price of a piece of length $i+1$
(1-indexed), cut the rod into pieces (each possible length usable
\textbf{any number of times}, since multiple pieces of the same length
can come from different parts of the rod) to \textbf{maximize} total
price.}

\begin{patternbox}
This closing problem of the chapter is a deliberate capstone connecting
its two halves: Rod Cutting is \textbf{0/1 Knapsack's value-maximization
recurrence} (Section~\ref{sec:dp-01-knapsack}), but in the
\textbf{unbounded}-reuse setting (Section~\ref{sec:dp-coin-change}'s
pick-stays-at-the-same-index rule) -- the ``items'' are piece lengths
$1, \ldots, n$, their ``weights'' are their own lengths, their
``values'' are their prices, and the ``capacity'' is the rod's total
length $n$.
\end{patternbox}

\subsection*{Deriving the recurrence}

Let $f(\textit{index}, \textit{length})$ be the maximum price obtainable
from a rod of $\textit{length}$, using piece lengths $1..\textit{index}+1$
(0-indexed $\textit{index}$ corresponds to piece length
$\textit{index}+1$). Either exclude this piece length entirely
($f(\textit{index}-1, \textit{length})$), or cut off one piece of this
length (gaining $\textit{price}[\textit{index}]$, then recursing on the
\emph{same} index since the same length remains available to cut again
from what's left):
\begin{gather*}
  f(\textit{index}, \textit{length}) = \max\Big(
  f(\textit{index}{-}1, \textit{length}),\ \textit{take} \Big), \\
  \textit{take} = (\textit{index}{+}1) \le \textit{length}\ ?\
  \textit{price}[\textit{index}] + f(\textit{index}, \textit{length}-(\textit{index}{+}1))
  : -\infty.
\end{gather*}
Base case: $f(0, \textit{length}) = \textit{length} \times
\textit{price}[0]$ (using only length-$1$ pieces, cut exactly
$\textit{length}$ of them).

\subsection*{Approach 1: Recursion}

\begin{lstlisting}
#include <bits/stdc++.h>
using namespace std;

int rodCuttingRecursive(int index, int length, vector<int>& price) {
    if (index == 0) return length * price[0];

    int notTake = rodCuttingRecursive(index - 1, length, price);
    int take = INT_MIN;
    int pieceLen = index + 1;
    if (pieceLen <= length) {
        take = price[index] + rodCuttingRecursive(index, length - pieceLen, price);
    }
    return max(notTake, take);
}

int main() {
    vector<int> price = {2, 5, 7, 8, 10};
    int n = price.size();
    cout << rodCuttingRecursive(n - 1, n, price) << endl; // 12
    return 0;
}
\end{lstlisting}

\begin{complexitybox}[Approach 1 Complexity]
\textbf{Time.} Exponential -- the same unbounded-reuse blowup as Coin
Change.
\textbf{Space.} $O(n)$ recursion stack.
\end{complexitybox}

\subsection*{Approach 2: Memoization}

\begin{lstlisting}
#include <bits/stdc++.h>
using namespace std;

int rodCuttingMemo(int index, int length, vector<int>& price,
                    vector<vector<int>>& memo) {
    if (index == 0) return length * price[0];
    if (memo[index][length] != -1) return memo[index][length];

    int notTake = rodCuttingMemo(index - 1, length, price, memo);
    int take = INT_MIN;
    int pieceLen = index + 1;
    if (pieceLen <= length) {
        take = price[index] + rodCuttingMemo(index, length - pieceLen, price, memo);
    }
    return memo[index][length] = max(notTake, take);
}

int main() {
    vector<int> price = {2, 5, 7, 8, 10};
    int n = price.size();
    vector<vector<int>> memo(n, vector<int>(n + 1, -1));
    cout << rodCuttingMemo(n - 1, n, price, memo) << endl; // 12
    return 0;
}
\end{lstlisting}

\begin{complexitybox}[Approach 2 Complexity]
\textbf{Time.} $O(n^2)$.
\textbf{Space.} $O(n^2)$ memo $+$ $O(n)$ recursion stack.
\end{complexitybox}

\subsection*{Approach 3: Tabulation}

\begin{lstlisting}
#include <bits/stdc++.h>
using namespace std;

int rodCuttingTabulation(vector<int>& price) {
    int n = price.size();
    vector<vector<int>> dp(n, vector<int>(n + 1, 0));

    for (int len = 0; len <= n; len++) dp[0][len] = len * price[0];

    for (int index = 1; index < n; index++) {
        for (int len = 0; len <= n; len++) {
            int notTake = dp[index - 1][len];
            int take = INT_MIN;
            int pieceLen = index + 1;
            if (pieceLen <= len) {
                take = price[index] + dp[index][len - pieceLen];
            }
            dp[index][len] = max(notTake, take);
        }
    }
    return dp[n - 1][n];
}

int main() {
    vector<int> price = {2, 5, 7, 8, 10};
    cout << rodCuttingTabulation(price) << endl; // 12
    return 0;
}
\end{lstlisting}

\subsection*{Dry run of the tabulation table}

\dryrun{Take $\textit{price} = [2,5,7,8,10]$ (rod length $n=5$).}

The optimal cut for a rod of length $5$ with these prices is two pieces
of length $2$ (price $5$ each) and one of length $1$ (price $2$):
$5+5+2=12$. Tracing the table: row $0$ (only length-$1$ pieces) gives
$\textit{dp}[0][5]=5\times2=10$. Row $1$ (lengths $1,2$) considers
mixing in length-$2$ pieces: $\textit{dp}[1][5] =
\max(\textit{dp}[0][5],\ \textit{price}[1]+\textit{dp}[1][3]) =
\max(10,\ 5+\textit{dp}[1][3])$; unrolling $\textit{dp}[1][3] =
\max(\textit{dp}[0][3],\ 5+\textit{dp}[1][1]) = \max(6, 5+2)=7$, giving
$\textit{dp}[1][5] = \max(10, 5+7)=12$. Later rows (introducing longer,
less efficient-per-length pieces) do not improve on this, so
$\textit{dp}[4][5]=12$.

\begin{complexitybox}[Approach 3 Complexity]
\textbf{Time.} $O(n^2)$.
\textbf{Space.} $O(n^2)$ table, no recursion stack.
\end{complexitybox}

\subsection*{Approach 4: Space Optimization}

\begin{lstlisting}
#include <bits/stdc++.h>
using namespace std;

int rodCuttingSpaceOptimized(vector<int>& price) {
    int n = price.size();
    vector<int> prev(n + 1, 0);

    for (int len = 0; len <= n; len++) prev[len] = len * price[0];

    for (int index = 1; index < n; index++) {
        vector<int> current(n + 1, 0);
        for (int len = 0; len <= n; len++) {
            int notTake = prev[len];
            int take = INT_MIN;
            int pieceLen = index + 1;
            if (pieceLen <= len) {
                take = price[index] + current[len - pieceLen]; // same row
            }
            current[len] = max(notTake, take);
        }
        prev = current;
    }
    return prev[n];
}

int main() {
    vector<int> price = {2, 5, 7, 8, 10};
    cout << rodCuttingSpaceOptimized(price) << endl; // 12
    return 0;
}
\end{lstlisting}

\begin{complexitybox}[Approach 4 Complexity]
\textbf{Time.} $O(n^2)$.
\textbf{Space.} $O(n)$ for the rolling row.
\end{complexitybox}

\begin{takeawaybox}
Rod Cutting closes this chapter by making its organizing idea explicit
one final time: every problem here is the same $(\textit{index},
\textit{target})$ table, varying along exactly two independent axes --
\emph{what the combination operator is} (OR, sum, or max/min) and
\emph{whether the pick branch reuses the current index or advances past
it}. Four operators (OR, sum, min, max) times two reuse rules (0/1,
unbounded) accounts for essentially every problem in this chapter, Rod
Cutting included.
\end{takeawaybox}

